\chapter{Introdução}
\label{chap:introducao}

\epigraph{\itshape ``A ciência avança um funeral de cada vez --- não porque os cientistas
mudam de ideia, mas porque as práticas experimentais que permitiam verificar isso
finalmente chegam ao alcance de todos.''}{\normalfont --- adaptado de Max Planck}

\section{Contexto e Motivação}
\label{sec:motivacao}

A robótica autônoma atravessa um momento singular em sua história. A convergência de
plataformas de software maduras --- especialmente o ROS~2 \cite{macenski2022ros2}
--- com hardware acessível e simuladores de alta fidelidade tornou possível, pela
primeira vez, que grupos de pesquisa sem grandes laboratórios físicos conduzam
experimentos de navegação com rigor comparável ao de instituições bem equipadas.
O que antes exigia uma sala climatizada, técnicos especializados e robôs industriais
de alto custo hoje pode ser realizado com um TurtleBot4, um computador convencional
e software inteiramente open-source.

Esse nivelamento de acesso, porém, não veio acompanhado de um nivelamento
metodológico equivalente. Enquanto os componentes técnicos se tornaram mais acessíveis,
os procedimentos que tornam um experimento reprodutível permaneceram informais,
específicos por equipe, raramente documentados e quase nunca transferíveis. A questão
não é mais \textit{se} um pesquisador consegue fazer um robô navegar --- é se consegue
fazer isso de uma forma que outro pesquisador possa verificar.

A comunidade de robótica reconhece essa lacuna metodológica há mais de uma
década \cite{bonsignorio2015}. Esta dissertação distingue, desde este parágrafo
em diante, dois conceitos que a literatura de ciência experimental frequentemente
trata como sinônimos, mas que não o são \cite{acm2020badging}. \textbf{Repetibilidade} é a capacidade do
mesmo sistema, no mesmo ambiente, sob as mesmas condições, de produzir
resultados equivalentes em execuções sucessivas. \textbf{Reprodutibilidade}, em
sentido mais forte, é a capacidade de um investigador independente, em outro
ambiente, obter resultados equivalentes a partir do mesmo artefato de pesquisa
--- código, dados e protocolo publicados \cite{amigoni2010}. Repetibilidade é
condição necessária, mas não suficiente, para reprodutibilidade: um sistema
pode ser perfeitamente repetível em um único laboratório e, ainda assim, nunca
ter sido reproduzido por ninguém fora dele.

Na prática, em robótica móvel, garantir repetibilidade significa assegurar que
o robô: (i) parta da mesma pose inicial em todas as execuções; (ii) percorra o
mesmo trajeto usando a mesma estratégia de navegação; (iii) tenha os dados de
todos os sensores coletados de forma síncrona e com a mesma configuração; e
(iv) tenha as trajetórias resultantes comparadas por um método que não
introduza vieses de análise. Cada um desses requisitos, isoladamente, é
simples. Satisfazê-los simultaneamente, de forma automática e reutilizável, é
onde a literatura apresenta uma lacuna persistente --- e é precisamente
repetibilidade, nesse sentido preciso, o que esta dissertação avalia
empiricamente (Capítulos~\ref{chap:avaliacao} e~\ref{chap:resultados}). O
framework proposto busca \emph{favorecer} reprodutibilidade --- por meio de
código aberto, protocolo documentado e artefatos de saída estruturados
(Seção~\ref{sec:rel_reprodutibilidade}) --- mas essa dissertação não conduziu
nem observou replicação independente por outro pesquisador ou laboratório;
a Seção~\ref{sec:repro_vs_repet} retoma essa distinção à luz dos resultados
obtidos.

Essa lacuna tem custos concretos. Quando um artigo reporta que ``o robô navegou com
erro médio de 8\,cm'', sem especificar se essa é a média de 3 ou 30 execuções, se o
desvio padrão foi medido, ou se a comparação com trabalhos anteriores usa a mesma
métrica, a contribuição científica se torna não verificável. A acumulação de
resultados não reprodutíveis na literatura não é apenas um problema acadêmico
--- é um obstáculo para a transferência de tecnologia para aplicações industriais
que dependem de garantias comportamentais precisas.

\section{O Problema Central: Crise de Orquestração Experimental}
\label{sec:problema}

Esta dissertação identifica e nomeia um problema específico que perpassa a robótica
experimental contemporânea: a \textbf{crise de orquestração experimental}. Trata-se
da dificuldade estrutural em coordenar, de forma automática e repetível, as múltiplas
camadas de um experimento de navegação autônoma --- inicialização do ambiente de
simulação, ativação da pilha de navegação com tempo de espera correto, definição e
execução de rotas, coleta síncrona de múltiplos sensores com políticas de QoS adequadas,
e análise quantitativa comparável entre execuções.

Cada uma dessas camadas possui sua própria curva de aprendizado. O modelo de lifecycle
do ROS~2 \cite{macenski2022ros2} exige que os nós de navegação sejam ativados na
sequência correta antes que qualquer comando seja enviado. As políticas de QoS de
tópicos como \texttt{/\allowbreak scan} (BEST\_EFFORT), \texttt{/\allowbreak pose} (RELIABLE) e
\texttt{/\allowbreak amcl\_\allowbreak pose} (TRANSIENT\_LOCAL) determinam se os dados serão efetivamente
gravados nos bags ou silenciosamente descartados. O protocolo de ações do Nav2
\cite{macenski2020marathon} opera de forma assíncrona, e uma chamada prematura ao
servidor de ações resulta em timeout sem mensagem de erro clara. O SLAM Toolbox
\cite{macenski2021slam} precisa de scans suficientes antes de publicar uma estimativa
confiável de mapa, e o intervalo correto varia com o ambiente.

Nenhum desses problemas é intratável. Todos são resolvidos por equipes experientes,
individualmente, a cada novo experimento. O problema é que essa solução acontece
de forma invisível, em scripts não publicados, em notebooks Jupyter apagados ao fim
do projeto, em conversas de corredor que nunca entram na seção de metodologia do artigo.
O conhecimento que torna um experimento bem-feito permanece tácito e não transferível.

\section{O Papel da Robótica de Frota nos Experimentos Modernos}
\label{sec:frota}

A tendência recente em robótica de serviço e logística autônoma é operar não um
robô isolado, mas frotas de agentes coordenados \cite{macenski2022ros2}. Em um
testbed experimental, isso cria possibilidades metodológicas que um único robô
não oferece: enquanto uma Unidade Móvel sob Teste (MUUT) percorre sua rota, uma
Unidade Fixa sob Teste (FUUT) pode observar a mesma cena de um ângulo externo,
coletando dados independentes que funcionam como fonte de verdade alternativa.

Essa dualidade de perspectivas --- endógena (o que o robô vê de si mesmo) e exógena
(o que um observador fixo vê do robô) --- enriquece substancialmente a análise de
repetibilidade. Pequenas variações de trajetória que seriam invisíveis à odometria
interna do MUUT podem ser capturadas pelos scans do FUUT como variações na posição
relativa do robô ao passar em frente ao sensor.

O framework proposto nesta dissertação formaliza esse design experimental com uma
arquitetura de papéis explícita, garantindo que as regras de cada tipo de unidade
sejam verificadas em software --- não apenas documentadas no protocolo.

\section{Escopo e Restrições}
\label{sec:escopo}

Esta dissertação foca na navegação 2D de robôs diferenciais em ambientes fechados,
usando o TurtleBot4 Standard com RPLIDAR S2 como plataforma de referência. Todos
os experimentos de validação foram conduzidos em simulação Gazebo Harmonic, o que
garante condições perfeitamente controladas mas abstrai os fenômenos físicos de
superfícies reais, slip de roda e variação de calibração de sensores ao longo do
tempo.

As métricas adotadas --- RMSE, erro de endpoint e consistência temporal --- são
métricas espaciotemporais clássicas para avaliação de navegação \cite{amigoni2010},
mas não capturam aspectos como consumo energético, desgaste mecânico ou robustez
a falhas de comunicação. Essas dimensões constituem extensões naturais do trabalho,
identificadas na seção de trabalhos futuros.

\section{Objetivos}
\label{sec:objetivos}

O objetivo geral desta dissertação é projetar e implementar um framework
open-source de orquestração de experimentos de repetibilidade de navegação
autônoma para frotas de robôs móveis em ROS~2, e avaliar empiricamente sua
capacidade de produzir execuções repetíveis --- avaliação que, como detalhado
no Capítulo~\ref{chap:avaliacao}, foi conduzida com uma única unidade móvel; a
extensão dessa avaliação para operação simultânea de múltiplos robôs é tratada
como trabalho futuro (Seção~\ref{sec:futuros}).

Os objetivos específicos são:

\begin{enumerate}
    \item Definir formalmente uma arquitetura de papéis experimentais (MUUT/FUUT/SU)
    que separe responsabilidades entre unidades móveis e fixas de forma explícita e
    verificável em software.

    \item Implementar serviços ROS~2 que abstraiam as ações Nav2 e a coleta síncrona
    de sensores em uma interface de alto nível, sem expor ao pesquisador os detalhes
    de QoS, lifecycle ou protocolo de ações.

    \item Desenvolver um protocolo de gravação--reprodução (record--replay) com
    garantias de consistência entre execuções, incluindo retorno automático ao ponto
    de partida e gravação de poses efetivamente atingidas.

    \item Construir uma interface web que permita configuração, execução e visualização
    de experimentos com waypoints clicáveis sobre o mapa SLAM ao vivo, sem necessidade
    de linha de comando.

    \item Validar o framework em dois cenários de simulação e reportar métricas de
    repetibilidade (RMSE inter-replay, RMSE vs. baseline, erro de endpoint e
    consistência temporal) como baseline empírico para trabalhos futuros.
\end{enumerate}


\section{Perguntas de Pesquisa e Hipóteses}
\label{sec:perguntas_pesquisa}

A partir dos objetivos específicos listados na Seção~\ref{sec:objetivos}, esta
dissertação formaliza quatro perguntas de pesquisa (PP), cada uma associada a
uma hipótese de trabalho (H) verificada empiricamente nos capítulos de
Avaliação e Resultados (Capítulos~\ref{chap:avaliacao} e~\ref{chap:resultados}).

\begin{description}
    \item[PP1 --- Automação de protocolo.] É possível automatizar integralmente
    o protocolo de gravação--reprodução de rotas em uma frota heterogênea de
    robôs ROS~2, sem intervenção manual entre fases, preservando a semântica
    de papéis experimentais (MUUT/FUUT/SU)?
    \textbf{H1}: um conjunto reduzido de serviços ROS~2 de alto nível
    (Seção~\ref{sec:orquestrador}), combinado a um script orquestrador de
    linha de comando, é suficiente para expressar qualquer protocolo de
    record--replay sem código adicional por experimento.

    \item[PP2 --- Repetibilidade mensurável.] Qual é o grau de repetibilidade
    espacial (RMSE) e temporal (CV$_t$) de execuções de replay sucessivas sobre
    a mesma rota, sob controlador e SLAM padrão do Nav2/SLAM Toolbox?
    \textbf{H2}: em condições de simulação controlada, o RMSE inter-replay é
    inferior a 30\,cm e o CV$_t$ é inferior a 10\%. Com $N=10$ replays na
    campanha final, esta dissertação caracteriza a variabilidade observada no
    cenário avaliado, sem generalizar esse comportamento para todas as rotas,
    ambientes ou configurações do Nav2 (Seção~\ref{sec:mppi}).

    \item[PP3 --- Persistência da repetibilidade fora do regime estritamente
    controlado.] A baixa variabilidade observada sob o protocolo controlado
    (mesma ação de navegação, pose inicial fixa, $N=10$ na campanha final) se
    mantém quando as condições de execução são menos controladas?
    \textbf{H3}: espera-se que execuções fora do protocolo formal --- sem
    controle explícito de ação de navegação, pose inicial e retorno ao ponto
    de partida entre réplicas --- apresentem RMSE e variabilidade de duração
    substancialmente maiores que a campanha controlada final. Esta dissertação não
    hipotetiza, a priori, uma causa única para essa diferença; a Seção
    ~\ref{sec:res_exp_complementares} discute evidência exploratória
    compatível com H3, mas insuficiente para atribuir causalidade específica,
    e a Seção~\ref{sec:ameacas} trata essa limitação como ameaça à validade
    de conclusão.

    \item[PP4 --- Custo de engenharia para o pesquisador.] O framework reduz
    o esforço de engenharia necessário para configurar e executar um
    experimento de repetibilidade, em comparação com a composição manual de
    chamadas Nav2/rosbag2 descrita informalmente na literatura correlata
    (Capítulo~\ref{chap:relacionados})?
    \textbf{H4}: a interface web e a camada de serviços ROS~2
    (Seção~\ref{sec:backend_web}) eliminam a necessidade de escrever código
    Python ou lançar comandos \texttt{ros2} manualmente para configurar,
    executar e analisar um experimento completo.
\end{description}

Este trabalho não formaliza PP2--PP4 como hipóteses estatísticas testadas por
inferência formal com poder amostral pré-calculado. Embora a campanha final
tenha $N=10$ replays, ela permanece restrita a uma única rota curta, uma única
plataforma simulada e uma única configuração de navegação, limitação reconhecida
explicitamente na Seção~\ref{sec:ameacas}. As hipóteses acima são tratadas,
portanto, como proposições descritivas a serem corroboradas ou refutadas pela
evidência empírica coletada, no espírito de um estudo exploratório de engenharia
de software experimental \cite{amigoni2010}, não de um experimento controlado
com validação estatística confirmatória.

\section{Contribuições}
\label{sec:contribuicoes}

As contribuições desta dissertação podem ser organizadas em três dimensões:

\subsection{Contribuições de Engenharia}

\begin{itemize}
    \item Implementação de \textbf{buffers TF por robô} para suporte a SLAM multi-robô
    em ROS~2, resolvendo o problema de namespace de tópicos \texttt{/\allowbreak tf} em frotas.

    \item Integração de \textbf{executor multithreaded} no orquestrador para evitar
    deadlocks entre callbacks de serviço e discovery de action servers.

    \item Desenvolvimento de \textbf{perfis de QoS diferenciados} no coletor de dados,
    garantindo que tópicos BEST\_EFFORT (sensores) e RELIABLE/TRANSIENT\_LOCAL
    (localização) sejam corretamente subscritos.

    \item Patch em tempo de launch do \textbf{Nav2 YAML} para ajuste de parâmetros
    de tipo de \texttt{cmd\_\allowbreak vel}, resolvendo incompatibilidade entre versões de
    simulação e navegação.

    \item Backend REST em FastAPI que expõe serviços ROS~2 sem dependência de
    ROSBridge, reduzindo a configuração necessária para novos usuários.
\end{itemize}

\subsection{Contribuições Científicas}

\begin{itemize}
    \item Formalização de uma \textbf{arquitetura de papéis experimentais}
    (MUUT/FUUT/SU) para testbeds de robótica com frota, não identificada nos
    trabalhos analisados nesta dissertação (Capítulo~\ref{chap:relacionados}).

    \item Definição e implementação de um \textbf{pipeline end-to-end} para
    experimentos de repetibilidade de navegação, do waypoint ao relatório RMSE.

    \item Resultados empíricos de repetibilidade com Nav2 e SLAM Toolbox em
    simulação Gazebo Harmonic, servindo como baseline para trabalhos futuros
    com hardware real e outros algoritmos de controle.

    \item Evidência exploratória de que execuções fora do protocolo controlado
    de record--replay (Seção~\ref{sec:res_exp_complementares}) podem produzir
    métricas contaminadas por variação de pose inicial --- achado que motiva,
    mas não confirma, hipóteses específicas como drift cumulativo de localização
    do SLAM, e que é reportado nesta dissertação como questão em aberto, não
    como resultado conclusivo.
\end{itemize}

\subsection{Contribuições para a Comunidade}

\begin{itemize}
    \item Código-fonte open-source publicado com documentação suficiente para
    adoção por outros grupos de pesquisa.

    \item Interface web acessível que democratiza o acesso a experimentos de
    repetibilidade, reduzindo a barreira de entrada para pesquisadores sem
    formação específica em ROS~2.
\end{itemize}

\section{Organização da Dissertação}
\label{sec:organizacao}

O restante desta dissertação está organizado da seguinte forma.

O \textbf{Capítulo~\ref{chap:fundamentacao}} apresenta o embasamento teórico sobre
navegação autônoma, cinemática diferencial, SLAM, o ecossistema ROS~2 e Nav2,
e as métricas clássicas de avaliação de repetibilidade de trajetórias.

O \textbf{Capítulo~\ref{chap:relacionados}} discute trabalhos relacionados em
frameworks de navegação, avaliação de SLAM e metodologias de teste em robótica,
posicionando a proposta em relação ao estado da arte.

O \textbf{Capítulo~\ref{chap:arquitetura}} descreve a arquitetura do framework
proposto, incluindo a formalização dos papéis experimentais, o orquestrador de frota,
o coletor de dados e a interface web.

O \textbf{Capítulo~\ref{chap:metodologia}} detalha o protocolo record--replay,
as métricas adotadas e o procedimento de análise de repetibilidade.

O \textbf{Capítulo~\ref{chap:implementacao}} apresenta decisões de implementação,
integração com o ambiente de simulação Gazebo Harmonic e a interface web desenvolvida.

O \textbf{Capítulo~\ref{chap:avaliacao}} descreve o arranjo experimental, as
questões de avaliação e os critérios de validação adotados.

O \textbf{Capítulo~\ref{chap:resultados}} analisa os resultados obtidos nos dois
cenários de validação e discute suas implicações.

O \textbf{Capítulo~\ref{chap:conclusao}} conclui com um resumo das contribuições,
limitações identificadas e a agenda de trabalhos futuros.

O \textbf{Glossário} reúne os termos técnicos centrais utilizados ao longo da
dissertação com suas definições precisas.
